% r03_hyb_coer: proof of \Cref{rot:prop:BD}, part II: coercivity at a regular background, L^2 bounds, Combes--Thomas decay, the decay of C^{(r)}.

\paragraph{Coercivity and decay in $L^2$ form.}
The proof of \Cref{rot:prop:BD}(a) follows \cite[\S2]{Ba4} and \cite[\S5.3]{K2} with two changes: the lower bound on the quadratic form is proved directly at a regular background by a cover with axis cubes (so that no small-domain reduction is needed for the $L^2$ theory), and the local cells of the random walk are periodic cells fitted by both lattices (\cite[(5.55)]{K2}), on which the $A=0$ kernels are the twisted periodisations of the explicit kernels of \Cref{rot:lem:H,rot:lem:slice}. On the tori of the chain the periods lie in $10\Z^3$ in the unit scale of every level $r\le k$ ($P_J=50MJ\,5^{k-r}\Z^3$ and $R^{-1}P_J=10MJ\,5^{k-r}(5R^{-1})\Z^3$, \Cref{rot:lem:adm}(iii)), with sup-systole $s_\infty(P)\ge50$ for $r\le k-1$.

\begin{lemma}[coercivity at a regular background]\label{rot:hyb:lem:coer}
Let $\Omega$ be $\Lambda_X$ or $\Lambda_X/P$ with $P$ fitted by both lattices and $s_\infty(P)\ge10$ in unit scale. There are $c_0(a)>0$ and $\gamma_0(a)>0$ such that for every regular $A$ on $\Omega$ with $e_r^\beta\le c_0$ and every $\phi$,
\[
\ip{\phi}{(-\Delta_A+m_r^2+a_rP_X(A))\phi}\ \ge\ \gamma_0\|\phi\|^2 .
\]
Explicitly $\gamma_0=\frac1{2m_c}\min\{\frac18,\frac{a'}{64(8a'+1)}\}$ with $a'=a(1-5^{-2})\le a_r\le a$ and $m_c$ the multiplicity of the cover below; $m_c=8$ when $P\subset2\Z^3$, as on the tori of the chain.
\end{lemma}
\begin{proof}
\emph{Cover.} Let $\{\Delta_i\}$ be the axis cubes $y_i+[0,4)^3$ with $y_i\in\Lambda_c\defeq2\Z^3+P$ ($P=0$ on $\Lambda_X$); $\Lambda_c$ is $P$-invariant, so the family descends to a finite family on the torus, and a cube does not wrap since every non-zero period has $|\gamma|_\infty\ge10$. Every point, every bond and every block lies in at most $m_c$ cubes, $m_c=8$ if $P\subset2\Z^3$ and $m_c\le64$ in general ($\Lambda_c\subset\Z^3$). Every $\Delta_i$ contains a full block: the covering radius of $\Gamma_r=R^{-1}\Z^3$ is $\frac{\sqrt3}2$, so some $z_i\in\Gamma_r$ has $|z_i-c_i|\le\frac{\sqrt3}2$ ($c_i$ the centre of $\Delta_i$), and $B^X_r(z_i)\subset\{|x-z_i|<\frac{\sqrt3}2\}\subset\{|x-c_i|<\sqrt3\}\subset\Delta_i$ because $\sqrt3<2$. Dropping the non-negative terms of the quadratic form that belong to bonds joining different cubes, to the mass and to the pins of the other blocks, and dividing by the multiplicity,
\[
\ip{\phi}{(-\Delta_A+m_r^2+a_rP_X(A))\phi}\ \ge\ \tfrac1{m_c}\sum_i\Big[\sum_{b\subset\Delta_i}\eta'^3|D_A\phi(b)|^2+a_r|\ip{w_{z_i}}{\phi}|^2\Big]=:\tfrac1{m_c}\sum_iq_i(\phi).
\]
\emph{Local gauge.} On the cube $\Delta_i$ (contractible) write $A=A_0^{(i)}+A'^{(i)}$ with $A_0^{(i)}\defeq A(c_i)$ constant and, by regularity along lattice paths of $\ell^1$-length $\le12$, $|A'^{(i)}|\le12\,Ce_r^{\beta-1}$ on $\Delta_i$. The gauge function $\lambda_i(x)=A_0^{(i)}\cdot x$ removes $A_0^{(i)}$ from both terms of $q_i$ (\Cref{rot:hyb:lem:alg}(ii),(iii), applied to the bonds and the block inside $\Delta_i$), so $q_i(\phi)=q_i^{A'}(\tilde\phi)$ with $\tilde\phi=e^{-\mathrm ie_r\lambda_i}\phi$, $|\tilde\phi|=|\phi|$.
\emph{Perturbation.} $D_{A'}\phi(b)=U(-A'_b)D_0\phi(b)+\eta'^{-1}(U(-A'_b)-1)\phi(b_-)$, hence $\|D_{A'}\phi\|_{\Delta_i}\ge\|D_0\phi\|_{\Delta_i}-e_r|A'|_\infty\|\phi\|_{\Delta_i}$ and $\|D_{A'}\phi\|^2\ge\frac34\|D_0\phi\|^2-3e_r^2|A'|_\infty^2\|\phi\|^2$. For the pin, $w_z(A')-w_z(0)$ has sup norm $\le e_r|A'(\Gamma)|\le2.1\,e_r|A'|_\infty$ (\Cref{rot:hyb:lem:len}(i)), so $\big||\ip{w_z(A')}{\phi}|^2-|\ip{w_z(0)}{\phi}|^2\big|\le(4.2e_r|A'|_\infty+4.41e_r^2|A'|_\infty^2)\|\phi\|_B^2$.
\emph{The free cube.} With $\psi\defeq\tilde\phi-c$ on $\Delta_i$, $c$ the mean, $\sum_{b\subset\Delta_i}\eta'^3|D_0\tilde\phi(b)|^2\ge\lambda_1\|\psi\|^2_{\Delta_i}$ with $\lambda_1=4\eta'^{-2}\sin^2(\pi\eta'/8)\ge\frac14$ (the first Neumann eigenvalue of the cube of $4/\eta'$ points per side; $\sin t\ge\frac2\pi t$), and, since $|B_i|\eta'^3=1$, $|\ip{w_{z_i}}{\tilde\phi}|^2=|c+\ip{\mathbf 1_{B_i}}{\psi}|^2\ge(1-t)|c|^2-(\frac1t-1)\|\psi\|^2_{B_i}$ for $0<t<1$. Choosing $\frac1t-1=\frac{\lambda_1}{2a_r}$,
\[
\lambda_1\|\psi\|^2+a_r|\ip{w_{z_i}}{\tilde\phi}|^2\ \ge\ \tfrac{\lambda_1}2\|\psi\|^2+\tfrac{a_r\lambda_1}{2a_r+\lambda_1}|c|^2\ \ge\ \gamma_1\|\tilde\phi\|^2_{\Delta_i},\qquad
\gamma_1\defeq\min\Big\{\tfrac{\lambda_1}2,\tfrac{a_r\lambda_1}{64(2a_r+\lambda_1)}\Big\},
\]
using $\|\tilde\phi\|^2_{\Delta_i}=\|\psi\|^2+64|c|^2$. \emph{Assembly.} $q_i(\phi)\ge\frac34\gamma_1\|\phi\|^2_{\Delta_i}-(3e_r^2|A'|^2+a_r(4.2e_r|A'|+4.41e_r^2|A'|^2))\|\phi\|^2_{\Delta_i}\ge\frac12\gamma_1\|\phi\|^2_{\Delta_i}$ once $e_r^\beta\le c_0(a)$, and $\sum_i\|\phi\|^2_{\Delta_i}\ge\|\phi\|^2$. With $\lambda_1\ge\frac14$, $\gamma_1\ge\min\{\frac18,\frac{a'}{64(8a'+1)}\}$.
\end{proof}

This is \cite[(5.59)--(5.60)]{K2} (``cover $\square$ by cubes $\{\Delta_i\}$ on the $\eta$-lattice with sides parallel to the coordinate planes \dots\ each cube contains fully at least one block'') written out, with the local gauge step that \cite[p.~345]{K2} leaves to ``$A=B+A'$ with $B$ constant and $A'$ small'' made explicit; it is also \cite[(1.8)]{Ba4} for the hybrid flow. No Poincar\'e inequality on the digitised rotated cube is needed: only axis cubes of the fine lattice enter.

\begin{corollary}[$L^2$ bounds]\label{rot:hyb:cor:L2}
Under the hypotheses of \Cref{rot:hyb:lem:coer}, $T\defeq-\Delta_A+m_r^2+a_rP_X(A)\ge\gamma_0$, so $\|G^X_r(A)\|\le\gamma_0^{-1}$, $\|D_AG^X_r(A)^{1/2}\|\le1$, $\|G^X_r(A)^{1/2}D_A^*\|\le1$, $\|D_AG^X_r(A)D_A^*\|\le1$ and $0\le G^X_r(A)\le\gamma_0^{-1}$ (\cite[Lemma~2.1]{Ba4}, \cite[Lemma~5.7]{K2}). The same holds, on $\Lambda_X$ and on a torus with $s_\infty(P)\ge10L$ in unit scale (true on the tori of the chain for $r\le k-1$), for the operator with the level-$(r+1)$ blocks, $-\Delta_A+m_r^2+a_{r+1}L^{-2}P_{r+1}(A)$ (the pin at scale $L$), with $\gamma_0$ replaced by $\gamma_0L^{-2}$.
\end{corollary}
\begin{proof}
$D_A^*D_A\le T$. For the $(r+1)$-blocks repeat the proof with cubes of side $4L$, corners in $2L\Z^3+P$ (multiplicity $8$ if $P\subset10\Z^3$, as for $P_J$ and $R^{-1}P_J$; periods with $|\gamma|_\infty\ge10L$) and blocks of $L^3$ unit cells (\cite[p.~588]{Ba4}: ``replace $\min\{\pi^2,a_k\}$ by $\min\{\pi^2L^{-2},a_{k+1}L^{-2}\}$'').
\end{proof}

\begin{proposition}[Combes--Thomas decay]\label{rot:hyb:prop:CT}
Under the hypotheses of \Cref{rot:hyb:lem:coer} there are $\theta_0,c_1>0$ depending on $a$ only such that for all $f,f'$ supported in unit cubes,
\[
|\ip f{G^X_r(A)f'}|,\ |\ip f{D_AG^X_r(A)f'}|,\ |\ip f{G^X_r(A)D_A^*f'}|\ \le\ c_1e^{-\theta_0\dist(\supp f,\supp f')}\|f\|_2\|f'\|_2 ,
\]
$\dist$ the quotient distance of $\Omega$. Hence (\cite[(3.6)]{K1}) $|A^X(A)(z,w)|=|a_r\delta_{zw}-a_r^2\ip{w_z}{G^X_r(A)w_w}|\le Ce^{-\theta_0|z-w|}$. The same holds for the operator with the $(r+1)$-pin of \Cref{rot:hyb:cor:L2}, with $\theta_0,c_1$ depending on $a$ and $L$.
\end{proposition}
\begin{proof}
For a weight $\psi$ with $|\psi(x)-\psi(x')|\le|x-x'|$ and $0\le t\le\theta_0$ put $T_\psi\defeq e^{t\psi}Te^{-t\psi}$. On $\Lambda_X$ one may take the linear weight $\psi=\hat n\cdot x$; on a torus $\vartheta\cdot x$ is not a function, and $\psi=\min\{\dist(\cdot,\supp f'),R\}$ is used, with the quotient distance; the estimates below use only the Lipschitz bound on bonds and inside blocks. The hopping part changes by a form-bounded perturbation: $e^{t\psi}D_Ae^{-t\psi}=D_A+B_t$, with $B_t$ the bond multiplication by $\eta'^{-1}(e^{-t(\psi(b_+)-\psi(b_-))}-1)$ composed with the covariant shift (bond functions weighted at $b_-$), of norm $\|B_t\|\le\sqrt3\,te^{t\eta'}$ (every site is the end $b_+$ of three bonds); with $e^{t\psi}D_A^*e^{-t\psi}=(D_A+B_{-t})^*$ this gives $|\ip\phi{(T_\psi-T)\phi'}|\le Ct(\|D_A\phi\|\|\phi'\|+\|\phi\|\|D_A\phi'\|)+Ct^2\|\phi\|\|\phi'\|$ for $t\eta'\le1$. The pin changes by $\sum_zw_z(x)\overline{w_z(y)}(e^{t(\psi(x)-\psi(y))}-1)$ with $|\psi(x)-\psi(y)|\le|x-y|\le\sqrt3$ inside a block, an operator of norm $\le2\sqrt3\,te^{\sqrt3t}$ (one rank-one piece per block, orthogonal supports). Hence $\|T^{-1/2}(T_\psi-T)T^{-1/2}\|\le Ct\gamma_0^{-1/2}+Ct^2\gamma_0^{-1}\le\frac12$ for $t\le\theta_0(a)$, using $\|D_AT^{-1/2}\|\le1$ (\Cref{rot:hyb:cor:L2}); so $T_\psi=T^{1/2}(1+T^{-1/2}(T_\psi-T)T^{-1/2})T^{1/2}$ is invertible with $\|T_\psi^{-1}\|\le2\gamma_0^{-1}$. Then $\ip f{G^X_rf'}=\ip{e^{-t\psi}f}{T_\psi^{-1}e^{t\psi}f'}$, since $G^X_r=e^{-t\psi}T_\psi^{-1}e^{t\psi}$; choose $t=\theta_0$ and $\psi=\min\{\dist(\cdot,\supp f'),R\}$ with $R\ge\dist(\supp f,\supp f')$ (or, on $\Lambda_X$, $\psi=\hat n\cdot x$ with $\hat n$ the unit vector from $\supp f'$ to $\supp f$), and use $\sup_{\supp f}e^{-t\psi}\sup_{\supp f'}e^{t\psi}=e^{-\theta_0\dist(\supp f,\supp f')}$ for the distance weight ($\psi=0$ on $\supp f'$), $\le e^{-\theta_0(\dist-2\sqrt3)}$ for the linear one. The derivative versions use $\|D_AT_\psi^{-1}\|,\|T_\psi^{-1}D_A^*\|\le C$ (the same algebra). For $A^X$: $w_z,w_w$ have unit norm and supports in blocks of diameter $\le\sqrt3$. The same proof with the $(r+1)$-pin (blocks of diameter $\le\sqrt3L$, \Cref{rot:hyb:cor:L2}) gives the bounds for that operator.
\end{proof}

\begin{corollary}[{\cite[(3.6)]{K1} for $C^{(r)}$}]\label{rot:hyb:cor:Ck}
$|C^{(r),X}(A)(z,w)|\le Ce^{-\delta_0d(z,w)}$ in the quotient distance of the domain. For two domains $\Omega_1$, $\Omega_2$ as in \Cref{rot:hyb:lem:domdiff} (a common fundamental region $F$, a union of $(r+1)$-blocks, on which the backgrounds agree up to a collar of width $W_0$) and $y,y'\in\Gamma_r\cap F$,
\[
|C^{(r)}(\Omega_1,A_1;y,y')-C^{(r)}(\Omega_2,A_2;y,y')|\le Ce^{-\delta(d_\wedge(y,y')+\dist(y,\partial F)+\dist(y',\partial F))},
\]
$d_\wedge$ the smaller of the two quotient distances; the same for $\Delta^{(r)}=A^X$.
\end{corollary}
\begin{proof}
\cite[\S3]{Ba4} writes $C^{(r)}(\Omega,A;y,y')$ through the Gaussian integral (3.1), factorises it over the next-level blocks after the gauge transformation ``$\psi(y)\to U(A(\Gamma_{y,z}))\psi(y)$'' (p.~587), and arrives at the identity (3.8): $C^{(r)}$ is a bounded combination of $\ip{q_r(y)}{G_r(\Omega,\Lambda,A)q_r(y')}$ and $\ip{q_{r+1}(y)}{\cdot\,q_{r+1}(y')}$, with $q_r(y)$, $q_{r+1}(y)$ the block modes (3.7) of unit and $L^{-d/2}$ norm, where $G_r(\Omega,\Lambda,A)$ is the fluctuation propagator pinned on the $(r+1)$-blocks (3.5). The two inputs of this algebra are that the blocks partition the lattice and that every contour runs from the block root to the point, both true for the hybrid flow (\Cref{rot:hyb:lem:alg}(i),(iv), \Cref{rot:lem:tile}); the decay of $G_r(\Omega,\Lambda,A)$ in the $L^2$-weighted form is \Cref{rot:hyb:prop:CT} for the $(r+1)$-pin. This is what \cite[pp.~345--346]{K2} says: ``To prove Proposition 5.6, we use the method presented in Chap.~3 of [Ba4]. This method requires only $L^2$-bounds on the operators $G_k(\Omega,A)$.'' The two-domain bound is not \cite[(3.9)--(3.10)]{Ba4}, which change the pin region inside one domain: on each domain (3.6) gives $C^{(r)}=V^*G_r(\Omega,\Omega^{(r)},A)V-(\beta/a_r)P+a_r^{-1}$ with $G_r(\Omega,\Omega^{(r)},A)$ the $(r+1)$-pinned operator (3.5), $V=Q_r^*(1-\beta P)$, $\beta=a_{r+1}L^{-2}/a_r$ ($Q_r=Q_X(A)$, $Q=Q^{(1)}(A)$, $P=Q^*Q$), whose local terms and block functions agree on the common region, so that $\delta C^{(r)}(y,y')=\ip{v(y)}{\delta G\,v(y')}$ there; $\delta G$ for the $(r+1)$-pin obeys the weighted bound \eqref{rot:hyb:eq:B1} of the geometric resolvent identity $\ip{f'}{(G_2-G_1)f}=\ip{f'}{G_1[H_2,\chi]G_2f}$ with \Cref{rot:hyb:prop:CT} on each side (\Cref{rot:hyb:lem:domdiff}).
\end{proof}
